
% ============================================================================
% PÁGINA DE TÍTULO
% ============================================================================
\begin{center}
  \vspace*{0.5cm}

  {\Large\textbf{\color{uv-dark}Declaración Responsable de Autoría y Uso de Inteligencia Artificial (IA) en el TFG}}

  \vspace{0.5cm}

  {\LARGE\color{uv-blue}ETSE-UV}

\end{center}

% ============================================================================
% SECCIÓN 1: HERRAMIENTAS DE IA UTILIZADAS
% ============================================================================
\seccion{Herramientas de IA Utilizadas}

Declaro que en este trabajo he utilizado las siguientes herramientas de IA (si procede):

\begin{enumerate}
  \item Herramienta y versión: \campo{}
  \item Herramienta y versión: \campo{}
  % Añade más herramientas si es el caso
\end{enumerate}

\vspace{0.5cm}

% ============================================================================
% SECCIÓN 2: DIMENSIONES DE USO DE LA IA
% ============================================================================
\seccion{1) Dimensiones de Uso de la IA}

\subtitulo{Instrucciones: Determina el grado de uso de la IA en las siguientes dimensiones. Marca una opción por fila.}

\noindent
\textit{Escala (5 puntos): \textbf{1} = No se ha utilizado IA, \textbf{2} = Uso muy puntual, \textbf{3} = Uso moderado,  \textbf{4} = Uso frecuente, \textbf{5} = Uso intensivo}

\begin{center}
  \small
  \setlength{\tabcolsep}{0.3cm}
  \begin{longtable}{|p{11cm}|c|c|c|c|c|}
    \hline
    \textbf{Dimensión} & \textbf{1} & \textbf{2} & \textbf{3} & \textbf{4} & \textbf{5} \\
    \hline
    \endhead
    \hline
    \endfoot

    Corrección lingüística: ortografía, gramática, puntuación y coherencia formal del texto &  &  &  &  & \\
    \hline

    Reescritura y mejora de estilo: reformulación para aumentar claridad, concisión y legibilidad manteniendo el significado &  &  &  & &  \\
    \hline

    Síntesis de contenido: elaboración de resúmenes a partir de textos extensos o complejos para identificar ideas clave &  &  &  &  &  \\
    \hline

    Apoyo en la localización de referencias: propuesta de posibles fuentes (a verificar y citar correctamente) &  &  &  & & \\
    \hline

    Planificación del trabajo: propuesta de esquemas, índices o guiones iniciales para estructurar el documento &  &  &  &  &  \\
    \hline

    Mejora de la argumentación y la estructura: detección de huecos, redundancias o saltos lógicos y propuesta de reorganización &  &  &  &  &  \\
    \hline

    Programación (si procede): apoyo en el desarrollo, depuración u optimización de fragmentos de código puntuales &  &  &  &  &  \\
    \hline

    Programación (si procede): uso generalizado de herramientas de generación de código de forma autónoma a partir de especificaciones &  & &  &  &  \\
    \hline

    Análisis de datos (si procede): ayuda exploratoria para limpieza preliminar, generación de hipótesis e identificación de patrones &  &  &  &  &  \\
    \hline

    Traducción o adaptación lingüística (si procede) &  &  &  &  &  \\
    \hline

    Soporte en material visual (si procede): ideas para figuras, tablas o gráficos y revisión de su explicación &  &  &  &  &  \\
    \hline

    Otros (especificar si procede)  &  &  &  &  &  \\
    \hline

  \end{longtable}
\end{center}


% ============================================================================
% SECCIÓN 3: COMPROMISOS DE USO RESPONSABLE
% ============================================================================
\seccion{2) Compromisos de Uso Responsable}

\subtitulo{Instrucciones: Determina si se ha realizado un uso responsable de la IA. Marca tu grado de acuerdo.}


\noindent
\textit{Escala (acuerdo): \textbf{1} = Totalmente en desacuerdo,  \textbf{2} = En desacuerdo, \textbf{3} = Neutral, \textbf{4} = De acuerdo, \textbf{5} = Totalmente de acuerdo}

\begin{center}
  \small
  \setlength{\tabcolsep}{0.3cm}
  \begin{longtable}{|p{11cm}|c|c|c|c|c|}
    \hline
    \textbf{Declaración} & \textbf{1} & \textbf{2} & \textbf{3} & \textbf{4} & \textbf{5} \\
    \hline
    \endhead
    \hline
    \endfoot

    Las ideas principales, interpretaciones y análisis críticos del trabajo son propios. & & & & &  \\
    \hline

    El uso de IA no ha sustituido mis competencias básicas de investigación, redacción y razonamiento académico. &  &  &  &  &  \\
    \hline

    He revisado, contrastado e integrado críticamente cualquier contenido generado con IA antes de incorporarlo. &  &  &  &  &  \\
    \hline

    He verificado la precisión y fiabilidad del contenido, incluyendo datos, afirmaciones y referencias bibliográficas. &  &  &  &  &  \\
    \hline

  \end{longtable}
\end{center}


% ============================================================================
% SECCIÓN 4: OBSERVACIONES
% ============================================================================

% Si no hay observaciones puedes comentar todas las líneas de esta sección
\seccion{Observaciones (Opcional)}

\flushleft

\vspace{0.5cm}

\vspace{0.5cm}

\newpage


% Estas dos secciones no se deben modificar


% ============================================================================
% SECCIÓN 5: DECLARACIÓN FINAL
% ============================================================================
\seccion{Declaración Final}


\noindent
\fbox{\begin{minipage}{0.98\textwidth}
    \small

  A través de esta declaración, afirmo que he descrito de manera transparente y fidedigna el uso de IA de este trabajo. Asumo la autoría y responsabilidad de la obra final y confirmo que no he introducido datos personales o confidenciales contraviniendo la normativa de protección de datos personales\footnotemark[1] y que he respetado las normas en materia de propiedad intelectual\footnotemark[2], así como las previsiones sobre integridad académica\footnotemark[3] de la Universitat de València.

\end{minipage}}

\vspace{0.5cm}
% ============================================================================
% SECCIÓN 6: FIRMA Y FECHA
% ============================================================================
\seccion{Firma y Autorización}

\vspace{1cm}

\noindent
Fecha: \campo{\hspace{4cm}}

\vspace*{1cm}

\noindent
Firma:
% Puedes poner la firma digital aquí. Si no tienes firma digital, entonces imprime esta página,
% firma, escanea e incluye esa página escaneada en lugar de esta en el documento final.

% ============================================================================
% NOTAS A PIE DE PÁGINA
% ============================================================================
\footnotetext[1]{\href{https://www.boe.es/boe_catalan/dias/2018/12/06/pdfs/BOE-A-2018-16673-C.pdf}{Normativa de protección de datos personales de la Universitat de València}}
\footnotetext[2]{\href{https://www.boe.es/buscar/act.php?id=BOE-A-1996-8930}{Normativa de propiedad intelectual de la Universitat de València}}
\footnotetext[3]{\href{https://www.uv.es/graus/normatives/Codi_convivencia_bones_practiques_UV.pdf}{Normativa sobre integridad académica de la Universitat de València}}

\newpage
